problem:  
Let \((M,\mathcal{D},g_{\mathcal{D}})\) be a compact, equiregular sub-Riemannian manifold of step \(2\), endowed with its Popp measure \(\mu_{\mathrm{Popp}}\).  
Fix a smooth complementary subbundle \(\mathcal{V}\) so that \(TM=\mathcal{D}\oplus\mathcal{V}\), and for each \(\lambda>0\), consider the “Riemannian approximation” metric \(g_\lambda = g_{\mathcal{D}} \oplus \lambda^2 g_{\mathcal{V}}\).  
Denote by \(\mathrm{Ric}_{g_\lambda}\) the Riemannian Ricci tensor of \(g_\lambda\) and by \(\mathrm{Ric}^{\mathcal{H}}_{\mathrm{ABR}}\) the horizontal Ricci-type tensor defined via the Agrachev–Barilari–Rizzi curvature operator of the sub-Riemannian Jacobi curve associated with \((M,\mathcal{D},g_\mathcal{D})\).  

Assume that \((M,g_\lambda)\) have uniformly bounded sectional curvatures \(|\mathrm{Sec}_{g_\lambda}| \le K\) independent of \(\lambda\) for horizontal directions and that the O’Neill-type tensors related to the decomposition \(TM=\mathcal{D}\oplus\mathcal{V}\) remain uniformly bounded.  

**Question:**  
Investigate whether there exists a scaling exponent \(\alpha>0\) such that the renormalized curvature tensors satisfy  
\[
\lim_{\lambda\to\infty} \lambda^{-\alpha}\, \mathrm{Ric}_{g_\lambda}(X,Y) = \mathrm{Ric}^{\mathcal{H}}_{\mathrm{ABR}}(X,Y),
\quad \forall X,Y\in\Gamma(\mathcal{D}),
\]
where the limit is understood in the \(C^0\)-topology on tensor fields.  
Further, determine whether this convergence—and the associated exponent \(\alpha\)—is *intrinsic*, i.e. independent of the chosen complement \(\mathcal{V}\). If not, characterize the geometric conditions on \((M,\mathcal{D})\) (e.g., contact, fat, or parallelizable distributions) under which such independence holds.  

Why is it a "good" problem:  
This formulation isolates a fundamental and concrete issue at the interface between Riemannian and sub-Riemannian curvature theories: the existence and intrinsic nature of a limiting curvature tensor. It asks whether the second-variation curvature encoded by Agrachev–Barilari–Rizzi’s Jacobi-curve framework can be viewed as a renormalized limit of the classical Ricci tensors along a Riemannian collapsing sequence.  
The problem is mathematically well-posed—it specifies all geometric objects involved, the scaling to test, and the precise sense of convergence—and yet remains entirely open even in classical settings such as contact metric structures.  
Its resolution would produce the first rigorous tensor-level bridge between Riemannian and sub-Riemannian curvature, unifying analytic and geometric approaches. Bridging these two curvature notions requires tools from geometric analysis, microlocal control theory, and collapse geometry. Its statement is simple, but solving it would fundamentally clarify how curvature persists or transforms in singular geometric limits, embodying the depth, cross-field reach, and conceptual simplicity of a genuinely *good* problem.